% 张三慧V2思考题25.1--25.12候选解析；原图25.41--25.43需原图独立重绘。
\section*{第25章\quad 几何光学：思考题解析}
\paragraph{思考题 25.1}
当光的波长 \(\lambda\) 远小于遮挡物、
孔径、折射率变化尺度 \(L\) 时，
波的衍射与干涉修正相对较小，
\[
\boxed{\lambda/L\ll1}
\]
便可采用光线和Fermat最短光程的几何光学。
严格说，在\emph{均匀介质}中
几何光学光线近似沿直线传播；
在非均匀介质中
\(\nabla n\ne0\) 则光线\emph{可弯曲}，
故“几何光学”不等同于所有介质中直线。

\paragraph{思考题 25.2}
若树叶空隙孔径小于太阳像的几何尺寸、
但又没有小到令衍射占主导，
每个小孔相当于针孔投影仪，
地面光斑近似是太阳光盘的倒像。
因此圆盘形太阳产生圆斑，
月亮遮挡后的太阳呈新月形则各孔产生
新月像。
\[
\boxed{\text{光斑形状反映有限角尺寸的太阳光源。}}
\]
非常靠近地面或孔径过大时，孔形状将逐渐显现；
孔径过小时也会出现明显的衍射模糊。

\paragraph{思考题 25.3}
设人的眼睛离地 \(h_e\)、头顶高 \(H\)，
脚底在地面0。
平面镜反射的入射角等于反射角，
连接眼睛与脚底的镜面反射点在
\(h_e/2\)，连接眼睛与头顶的反射点在
\((h_e+H)/2\)。
故最小镜长
\[
\boxed{L_{\min}=\frac{h_e+H}{2}-\frac{h_e}{2}
=\frac H2}.
\]
它与眼睛高度、人与镜子的水平距离
无关（理想竖直平面镜、单一观察点）。
镜应从 \(h_e/2\) 高度到
\((h_e+H)/2\) 高度布置。

\paragraph{思考题 25.4}
凸面镜使来自更大物方角度范围的光
反射到眼内，在相同镜片尺寸下获得更大视场。
近轴凸面镜的像为
\emph{正立、缩小、位于镜后的虚像}，
所以物体角尺寸变小；
观察者习惯用物体已知尺寸判断远近，
常误以为车比实际\emph{更远}。
\[
\boxed{\text{凸面镜扩大视野；镜中物体看起来更远。}}
\]
并非所有镜前位置的距离知觉误差都完全由
球面镜方程决定，实际车距需要其他线索判断。

\paragraph{思考题 25.5}
炎热路面使靠近地面的空气温度更高、
密度与折射率更低，
\(n(z)\) 自上而下逐渐减小。
近掠射方向的远处天空光在折射率梯度下
逐渐弯向上方进入人眼，
使观察者把光反向延长成近地面虚像，
好似水面倒映天空（\emph{下蜃景}）。
海上\emph{上蜃景}可由海水上方
较冷、较密的底层空气及逆温层产生，
折射率随高度变化方式与下蜃景不同，
使远方物体像出现在实际位置的上方。
\[
\boxed{\text{蜃景是分层空气中的连续折射而非真实水面。}}
\]
不同类蜃景的具体光线路径须由
温度—折射率剖面决定。

\paragraph{思考题 25.6}
在一般\emph{正常色散}的透明玻璃可见区，
折射率 \(n_{\rm violet}>n_{\rm red}\)。
光在介质中的\emph{相速率}
\[
v_{\rm ph}(\lambda)=\frac c{n(\lambda)}.
\]
因此
\[
\boxed{v_{\rm red}>v_{\rm violet}}.
\]
频率跨界面保持不变，进入玻璃后
波长随 \(n\) 缩短。
在异常色散波段不能不说明条件就照搬此结论。

\paragraph{思考题 25.7}
从折射率 \(n_1>n_2\) 的高折射率介质
向低折射率介质射入时，
Snell关系
\(n_1\sin i=n_2\sin r\)。
使折射角达90°的临界角为
\[
\boxed{i_c=\arcsin(n_2/n_1)}.
\]
若 \(i>i_c\)，理想无吸收界面没有
向低折射率侧传送\emph{远场法向能流}的传播折射波，
发生\emph{全反射}，但低折射率侧仍有
沿界面指数衰减的倏逝场。
从 \(n_1<n_2\) 侧射入时不存在同类全反射临界角。

\paragraph{思考题 25.8}
浸在水中的空气饼状凸透镜
相对折射率
\(n_{\rm rel}=n_{\rm air}/n_w=1/1.33<1\)。
薄透镜制造公式
\[
\frac1f=(n_{\rm rel}-1)
\left(\frac1{R_1}-\frac1{R_2}\right).
\]
双凸几何取 \(R_1=+0.30\,\mathrm m\)、
\(R_2=-0.30\,\mathrm m\)，得
\[
\boxed{f=\left[
\left(\frac1{1.33}-1\right)\frac2{.30}
\right]^{-1}\simeq-0.605\,\mathrm m}.
\]
所以
\(\boxed{\text{对水中的平行光为发散透镜}}\)。
忽略塑料袋膜厚和曲面非球面误差，
且采用题设“薄透镜”近似。

\paragraph{思考题 25.9}
用正实物距 \(s>0\)、出射侧实像距 \(s'>0\)
的薄透镜符号约定，\(1/s+1/s'=1/f\)，
\(m=-s'/s\)：
\[
\begin{array}{c|c|c}
\text{凸透镜实物位置}&\text{像位置}&\text{性质}\\\hline
s>2f&f<s'<2f&\text{实、倒立、缩小}\\
s=2f&s'=2f&\text{实、倒立、等大}\\
f<s<2f&s'>2f&\text{实、倒立、放大}\\
s=f&s'=\infty&\text{平行出射、有限屏无清晰像}\\
0<s<f&s'<0&\text{虚、正立、放大}
\end{array}
\]
凹透镜 \(f<0\) 对任意实物 \(s>0\)：
\(s'<0\)、像在入射侧、\(|s'|<|f|\)，
\(\boxed{m>0,\ |m|<1}\)，即\emph{虚、正立、缩小}。
球面镜的 \(s,s'\) 符号可用另一套镜面约定，
应先说明正方向，不能不转换就套同一公式。

\paragraph{思考题 25.10}
对凸透镜和真实物体，
\(s>f\) 时
\[
s'=\frac{fs}{s-f}>f,
\]
故任何\emph{有限实像}都不能落在
\(\boxed{0<s'\le f}\) 区域；
\(s=f\) 的像在无穷远而非有限光屏上。
凸透镜对 \(0<s<f\) 可形成同侧虚像，
\emph{不能}笼统说焦内物体“不能成像”。
凹透镜 \(f=-F<0\) 时
\[
s'=-\frac{Fs}{s+F},
\qquad 0<|s'|<F.
\]
所以对\emph{真实物体}，薄凹透镜
没有出射侧的实像，且虚像位于
透镜至物方焦点之间的有限区间；
如果入射的是汇聚光形成的\emph{虚物}，
凹透镜也可能产生实像，需要另行计算。

\paragraph{思考题 25.11}
实像由光束在成像面\emph{真实相交}，
可用放置在该位置的屏幕承接，
也可在适当的观察方向用眼直接看到；
虚像是出射光的反向延长线相交，
放置光屏不能在该\emph{虚像位置}接收到
真正汇聚的清晰像。
\[
\boxed{\text{实像可投屏；虚像不直接投屏。}}
\]
眼睛接受的总是真实传播到瞳孔的光；
眼脑把其来源方向外推到虚像位置。
“看得到”不等于“可以投到屏上”。

\paragraph{思考题 25.12}
对图25.43(a)凸透镜，
若物距超过焦距，像在出射侧某有限位置
真实形成；直接看实像时，
眼应处于\emph{出射光束已越过实像而重新发散}
并进入瞳孔的区域，即一般在实像之后。
若物距小于焦距，眼在透镜右侧接受
离透镜发散的光，追溯至左侧虚像。
图25.43(b)凹透镜对实物总形成左侧虚像，
眼应在其右侧出射光束覆盖区域。
作图使用：
(1) 平行主轴入射，凸透镜后经实焦点、
凹透镜后似从物方焦点发出；
(2) 过光心的傍轴光线近似不偏折。
\[
\boxed{\text{眼睛应处于真实出射光束经过的空间，
而非凭虚像的几何位置直接放眼。}}
\]
焦距和物距决定精确交点；
原图25.43逐条光线的TikZ重绘仍待复核。
